% title: 习题：概率极限理论
% renderer: btex
% legacy-title: 用户:Solution/ 习题: 概率极限理论
% source: https://bananaspace.org/legacy/wiki/%E7%94%A8%E6%88%B7:Solution/_%E4%B9%A0%E9%A2%98:_%E6%A6%82%E7%8E%87%E6%9E%81%E9%99%90%E7%90%86%E8%AE%BA
% summary: 从香蕉空间旧站迁移的原始 bTeX 页面。
\setdisplaytitle{习题：概率极限理论}
\sethtmltitle{习题：概率极限理论}

\section{2026春习题}
\subsection{Homework\quad1}

\begin{enumerate}
    \item  Are the following distributions infinitely divisible?

    \begin{enumerate}
        \item[(i)] \(\displaystyle \mathbb{P}(X = k) = a^k (1 + a)^{-k-1}, \ a > 0, \ k = 0, 1, 2, \dots;\)
        
        \item[(ii)] \(\displaystyle p_0 = \mathbb{P}(X = 0) = (1 + \alpha \lambda)^{-1/\alpha}, \ \alpha, \lambda > 0,\) and
        \[
        \mathbb{P}(X = k) = \left( \frac{\alpha \lambda}{1 + \alpha \lambda} \right)^k \frac{(1 + \alpha) \cdots (1 + (k - 1)\alpha)}{k!}, \quad k = 1, 2, \dots.
        \]
        
        \item[(iii)] The density function \(\displaystyle f(x) = \frac{1}{\pi (1 + x^2)}, \ x \in \mathbb{R}.\)
    \end{enumerate}

    \item  Let \(\varphi\) be a characteristic function. Suppose there is an integer sequence \(n_k\) satisfying \(\lim_{k \to \infty} n_k = \infty\) such that for every \(k \ge 1\), there exists a characteristic function \(\varphi_k\) satisfying \(\varphi = (\varphi_k)^{n_k}\). Show that \(\varphi\) is infinitely divisible.

    \item  Compute the Lévy-Khintchine representation for the following characteristic functions.

    \begin{enumerate}
        \item[(i)] \(\displaystyle \varphi(t) = (1 - it / \beta)^{-\alpha}, \ \alpha, \beta > 0;\)
        \item[(ii)] \(\displaystyle \varphi(t) = e^{-\theta |t|}, \ \theta > 0;\)
        \item[(iii)] \(\displaystyle \varphi(t) = (1 + t^2)^{-1}.\)
    \end{enumerate}

    \item  Let \(X\) and \(Y\) be independent, infinitely divisible random variables. If \(X + Y\) has Normal distribution, show that \(X\) and \(Y\) also have Normal distributions.

    \item  Let \(\{X_{nk}\}_{n \ge 1, 1 \le k \le n}\) be independent random variables having distributions
    \[
    \mathbb{P}(X_{nk} = k/n) = \mathbb{P}(X_{nk} = -k/n) = 1/n \quad \text{and} \quad \mathbb{P}(X_{nk} = 0) = 1 - 2/n.
    \]
    Let \(\displaystyle S_n := \sum_{k=1}^n X_{nk}\). Show that
    \begin{enumerate}
        \item[(i)] For any \(n \ge 1\), \(S_n\) is not infinitely divisible.
        \item[(ii)] There exists a non-Normal, infinitely divisible random variable \(S\) such that \(\displaystyle S_n \xrightarrow{d} S\).
    \end{enumerate}
\end{enumerate}